\section{Side reflections and weak mixing}\label{sec:conclusion}

The analytic argument is complete. We now apply Proposition~\ref{prop:rigidity}
to the lift of a billiard eigenfunction. Its local form will make the angle
condition decisive.

\begin{proposition}\label{prop:eigenfunctions}
Let \(Q\) be a nondegenerate Euclidean triangle with an angle irrational
relative to \(\pi\). If \(F\in L^2(Q\times S^1,\mu)\) satisfies
\[
 F\circ\Phi_t=e^{i\lambda t}F\quad\text{in }L^2
 \qquad(t\in\R)
\]
for some \(\lambda\in\R\), then \(F\) is constant almost everywhere.
If this constant is nonzero, then \(\lambda=0\).
\end{proposition}

\begin{proof}
First suppose \(F\) is bounded. Put the origin of the triangle coordinates
at a vertex whose angle \(\alpha\) satisfies \(\alpha/\pi\notin\mathbb Q\).
Let \(f\) be the lift of \(F\) to the double, with its orientation chosen
so that these coordinates are oriented on the first sheet.
Proposition~\ref{prop:rigidity} and the eigenfunction equation give
\[
 Xf=i\lambda f,\qquad Yf=0,
 \qquad\text{hence}\qquad \nabla f=i\lambda v f
\]
in every oriented flat chart, where \(v=(\cos\theta,\sin\theta)\).
Thus \(e^{-i\lambda x\cdot v}f(x,v)\) has zero spatial distributional
gradient. On a connected chart it is independent of position for almost
every direction. One can see this by taking its angular Fourier coefficients:
each has zero spatial distributional gradient and is therefore constant
on the connected base chart. In the first triangle interior this gives
\begin{equation}\label{eq:plane-wave}
 F(x,v)=c(v)e^{i\lambda x\cdot v}
\end{equation}
for a measurable bounded function \(c\) on \(S^1\).

Consider an open side. Reflection in its supporting line is the affine map
\(p\mapsto Rp+d\), where \(R\) is the linear orthogonal reflection and
\(d\in\R^2\). A crossing chart agrees with the first-sheet coordinates
and assigns coordinates \(x=Rp+d\) to the point \(p\) of the second sheet.
The chart direction \(v\) folds to the billiard direction \(Rv\) there.
Since \(R^T=R=R^{-1}\), Equation~\eqref{eq:plane-wave} gives
\[
 f(x,v)=F(p,Rv)=c(Rv)e^{i\lambda (x-d)\cdot v}
\]
on the second portion of the chart. The phase-removed function is independent
of position throughout the connected crossing chart, so its expressions on
the two open portions agree:
\begin{equation}\label{eq:side-relation}
 c(v)=c(Rv)e^{-i\lambda d\cdot v}
 \quad\text{for almost every }v\in S^1.
\end{equation}
This comparison uses open subsets of a regular chart and requires no
boundary trace for \(F\).

For the two sides through the origin, \(d=0\). Their reflections generate
a rotation through \(2\alpha\), under which \(c\) is invariant.
If \(\widehat c(k)\) is its \(k\)-th circle Fourier coefficient, this gives
\(\widehat c(k)=e^{2ik\alpha}\widehat c(k)\). Irrationality of
\(\alpha/\pi\) forces every coefficient except \(k=0\) to vanish.
Thus \(c\) is constant almost everywhere.

The third side has \(d\ne0\), since its supporting line does not pass
through the chosen vertex. If the constant \(c\) is nonzero,
\eqref{eq:side-relation} says \(e^{-i\lambda d\cdot v}=1\) almost
everywhere on the circle, and hence everywhere by continuity. The values
of \(d\cdot v\) fill the interval \([-|d|,|d|]\), so this is possible
only when \(\lambda=0\). Equation~\eqref{eq:plane-wave} now makes \(F\)
constant, proving the bounded case.

For an arbitrary \(L^2\) eigenfunction, use
\(\widetilde F=F/(1+|F|)\). The map \(z\mapsto z/(1+|z|)\) is bounded,
injective, and commutes with multiplication by unit complex scalars.
Consequently \(\widetilde F\) is a bounded eigenfunction with the same
frequency. The bounded case and injectivity give the assertion for \(F\).
\end{proof}

We recall the standard compact-operator proof of the spectral implication
needed to finish the argument; see, for example,
\cite[Theorem~6.1.4]{Peterson2015}. It applies
to the full continuous-time flow.

\begin{lemma}\label{lem:spectral}
Let \((\Omega,\mu)\) be a standard probability space and \(\Phi_t\) a
measure-preserving flow whose composition operators form a strongly
continuous unitary group on complex \(L^2(\mu)\). If the only flow
eigenfunctions are constants, then the product flow
\(\Phi_t\times\Phi_t\) is ergodic for \(\mu\otimes\mu\).
\end{lemma}

\begin{proof}
Write \(U_t h=h\circ\Phi_t\). The frequency-zero hypothesis first implies
ergodicity of \(\Phi_t\). If the product flow were not ergodic, there
would be a nonzero invariant kernel \(K\in L^2(\mu\otimes\mu)\) of mean
zero. Each one-variable marginal of \(K\) is invariant under the single
flow and thus constant. The zero total mean makes both marginals zero.

Define the complex-linear Hilbert--Schmidt operator
\[
 (Th)(x)=\int_\Omega K(x,y)h(y)\,d\mu(y).
\]
The zero marginals imply \(T1=0\) and that its range is contained in
\(L^2_0(\mu)=\{h:\int h\,d\mu=0\}\). Its restriction to this space
is nonzero, since a zero integral operator has zero \(L^2\) kernel.
Product invariance gives \(U_tT=TU_t\) by a change of variables in the
integral. Hence \(T^*T\), acting on \(L^2_0(\mu)\), is a nonzero
positive compact operator commuting with every \(U_t\).

A positive eigenvalue of \(T^*T\) has a nonzero finite-dimensional
eigenspace invariant under every \(U_t\). Commuting unitary operators
on that space have a common eigenvector \(h\ne0\). The associated
scalar character is continuous by strong continuity of \(U_t\), so
it has the form \(e^{i\lambda t}\) for some \(\lambda\in\R\).
This gives a flow eigenfunction in \(L^2_0(\mu)\), a contradiction.
\end{proof}

\begin{proof}[Proof of Theorem~\ref{thm:main}]
The billiard composition group is strongly continuous by the construction
in Section~\ref{sec:geometry}. Proposition~\ref{prop:eigenfunctions}
eliminates every nonconstant eigenfunction, so Lemma~\ref{lem:spectral}
proves ergodicity of the product flow.
\end{proof}
